\subsubsection{Initial Milestone}

\textbf{Regarding PPP: Have you estimated the project execution cost?}
\vspace{-5pt}

The project budget is €14,592.78, as detailed in Section \ref{sec:budget}. This corresponds to the full ``bill'' for the application, including all personnel, hardware amortization, software licensing, and contingency costs.

\vspace{2cm}

\textbf{Regarding Exploitation: How are currently solved economic issues (costs...) related to the problem that you want to address (state of the art)?}
\vspace{-5pt}

According to the problems stated in Section \ref{sec:problems-to-solve}, the current time-consuming, low-efficiency expense management process delivers minimal value while misallocating company personnel costs for both employees and finance staff.


\textbf{Regarding Exploitation: How will your solution improve economic issues (costs, etc.) with respect to other existing solutions?}
\vspace{-5pt}

The primary benefit of the proposed solution is the operational efficiency optimization. The current manual process requires approximately 225 hours per year, as estimated in Section \ref{sec:problems-to-solve}. The application is designed to reduce the time required for this activity by at least 75\%, in accordance with the non-functional requirement NFR6 from Section \ref{sec:non-functional-requirements}.

Based on this reduction and the actual personnel wages at FSP, the solution generates annual cost savings of approximately €5,062.5. Given the estimated development cost of €14,592.78, the return on investment (ROI) is projected to be achieved within 3 years. This demonstrates that the proposed solution is more economically viable than investing in external solutions, supporting its adoption as a sustainable long-term approach.

However, an initial budget misestimation could extend the ROI period. This is mitigated with the 27-hour built-in buffer (calculated in Section \ref{sec:temporal-planning}) and the MoSCoW prioritization strategy (employed in Section \ref{sec:functional-requirements}). As a result, a high-value product is most likely to be delivered, since \textit{Could have} requirements can be postponed in case of budget shortfalls. 

